\section{Discussion}
\label{sec:discussion}

Fusion yields no meaningful gain over the best single stream on any of the three datasets or on the full DIPSER release (at most $+0.010$ QWK, on DAiSEE). This agrees with accounts of multimodal failure~\cite{wang2020makes,peng2022balanced}: the gate collapses to one stream chosen by the optimisation rather than by how informative it is, and mean pooling weakens the single informative stream with uninformative ones. The components that help concern optimisation and supervision rather than fusion: the global link gives the head direct access to all inputs, and the soft target exploits the rater distribution and yields the best single model on the slice.

The full release changes the interpretation of the slice. On the slice, models appeared useful because they ranked 20 students; with 57 students, labels vary mainly between lessons, and the sensor signal is largely shared by the class. Neither richer or local moments, nor a face crop, nor pretrained face encoders, nor temporal or class context finds a student-specific signal, so the limitation lies in how little the per-student streams convey about labels that vary with the lesson, rather than in the encoder.  Each relaxation of the split introduces a leak: the same student a few seconds apart, classmates in the same lesson, and classmates' labels retrieved through class features. Accuracy hides all three, because it remains close to the majority rate.

\textbf{Recommendations.} (1) Report the best single-stream baseline under the same protocol as any fusion model. (2) With multiple raters, train on the vote distribution rather than the majority vote. (3) Interpret one-hot gates as stream selection, not fusion. (4) Do not transfer components across datasets without testing them, since the link helps on DIPSER but harms on DAiSEE. (5) For imbalanced ordinal labels, base conclusions on QWK and report accuracy together with the majority rate. (6) For claims about classrooms, hold out whole classes, since person-grouped folds still share lessons and, with class-level features, leak classmates' labels. (7) Report label-only references alongside the majority class; here, the other students of the same lesson give QWK \fcExpPrior. (8) When earlier labels of the same student are available, report label persistence~\cite{demasi2017meaningless}.

\subsection{Threats to validity}
\label{sec:threats}
On the slice, folds group student-sessions rather than persons; re-evaluating every run on person-disjoint windows leaves the main findings unchanged (no fusion above the best stream; gate versus core $p=0.79$; link 10 of 12 pairs, $p=0.021$; best stream $0.213$, best fusion $0.109$), but the soft-target gain is no longer significant ($p=0.078$). Nested selection removes selection bias but trains on 80\% of the fold subjects, so absolute values are not comparable with the earlier non-nested version~\cite{cawley2010overfitting}. Only complete windows are used, which removes 47\% of the slice and 19\% of the full release. The slice contains seven high-attention windows, the mean fusion gap is driven by one seed, and the gradient probes use a single seed and fold. The face stream of the slice experiments is the whole webcam frame; Section~\ref{sec:fc-crop} tests a true crop. The full-release attempts used fixed model families (MLP, gradient-boosted trees, GRU); the face encoders were frozen, so we conclude that moment features, pretrained embeddings and context do not help, not that an encoder fine-tuned on classroom data could not. Leave-one-class-out uses only three folds (20, 21 and 16 students), and its estimates are correspondingly noisy. The DAiSEE comparison uses only clips with all streams, and the EngageNet results rely on a subset of configurations and on frame-level samples.
